\section{Method}
\label{sec:method}

% TARGET: ~2.5 pages

% PLACEHOLDER — to be written

\subsection{Overview}
% High-level architecture figure reference
% Three-phase approach: (1) pretrain autoencoder on target, (2) build memory bank, (3) train translator

\subsection{Shared Encoder and Memory Bank}
\label{sec:memory_bank}
% Shared encoder E_\phi maps both source and target windows to latent space
% Memory bank M = {E_\phi(x_T^(j))} pre-encoded from target domain
% k-NN retrieval: for each source timestep, find k nearest target windows
% Naturally causal: retrieval uses only past context

\subsection{Cross-Attention Translation}
\label{sec:cross_attention}
% Cross-attention blocks attend over retrieved target examples
% n_cross_layers is task-dependent (3 for dense labels, 2 for sparse)
% Decoder produces translated output in input space

\subsection{Training Procedure}
\label{sec:training}
% Phase 1: Autoencoder pretrain on target data (reusable checkpoint)
% Phase 2: End-to-end with frozen downstream model
%   L = L_task + lambda_f * L_fidelity + lambda_r * L_range + lambda_s * L_smooth
% L_task: BCE/MSE through frozen model
% L_fidelity: MSE(translated, original) — prevents catastrophic divergence
% L_range: penalizes out-of-range values
% Memory bank refresh every K epochs

\subsection{Design Choices}
\label{sec:design_choices}
% Temporal attention mode: causal vs bidirectional (task-dependent)
% Feature gate: learnable per-feature importance weights
% Cross-domain normalization: affine renorm of source stats to target
